% ============================================================================
%  CHAPTER 2: LITERATURE SURVEY   (expanded, deep-dive version)
%  Replaces the entire Chapter 2 body in main.tex.
% ============================================================================

\chapter{Literature Survey}
\label{chap:literature}

The work that FinITR-AI builds on does not come from a single field. It sits at
the intersection of three communities that rarely talk to one another: the
people who build tax-filing software in India, the people who apply large
language models to financial questions, and the researchers who study how to
make language models reliable by giving them tools and grounding. Each of these
brings something the project needs, and each leaves a gap the project tries to
fill. The rule-based tools are accurate but cannot see the AIS and cannot reason
about anything outside their fixed flows. The general models reason fluently but
fail at the arithmetic and invent rules. The tool-augmentation literature offers
the right techniques --- deferring computation to tools, retrieving grounded
evidence, checking claims for faithfulness --- but supplies them as
domain-agnostic frameworks rather than as a finished, verified tax system. This
chapter examines all three in turn. For each category we describe the
representative systems in enough detail to make the comparison meaningful, and
then we state precisely where our design diverges and why.

\section{Existing Approaches -- Rule-Based and Commercial ITR Tools}
\label{sec:category-a}

The systems that the overwhelming majority of Indian taxpayers actually use are
deterministic and form-driven. They encode the provisions of the Income Tax Act
as explicit logic, walk the user through a guided sequence of questions, and
validate the assembled return against the department's schema. On the standard
salaried case they are reliable, and the arithmetic they produce is, for the
most part, correct. Their limitation is one of perspective rather than of
accuracy: they reason over data the user supplies, they do not interpret the AIS
into actionable findings, and they cannot explain or adapt to a case that falls
outside their predefined paths.

\subsection{The Government e-Filing Portal and Offline Utilities}
The Income Tax Department publishes both an online portal and downloadable
offline utilities (in JSON and Excel form) that constitute the reference
implementation of the rules \cite{ref1}. These tools pre-fill certain fields
from the AIS and Form 26AS and enforce schema-level validation on the final
return, which catches structural errors before submission. What they do not do
is advise. They assume the filer already knows which form applies, which
schedules to attach, and which regime serves them better. The AIS is presented
as a list of raw entries rather than as a reconciled comparison against the
return being prepared, so the burden of spotting a mismatch --- the very thing
that triggers a notice --- remains entirely with the user.

\subsection{ClearTax}
ClearTax is among the most widely used assisted-filing platforms and automates
much of the common salaried journey: importing Form 16, walking the user through
deduction entry, offering a regime comparison, and handling e-filing
\cite{ref2}. It is a polished, capable product within its design envelope. The
constraints relevant to this project are two. First, its reconciliation of the
AIS against self-reported income is shallow --- it surfaces AIS data but does
not treat the gap between AIS and the return as a quantified risk to be
explained. Second, its reasoning is a fixed rule engine, which means it handles
the cases its designers anticipated and degrades on combinations they did not,
without the ability to explain why a particular treatment was chosen.

\subsection{Quicko}
Quicko offers a comparable guided-filing experience and is notably strong on
capital-gains import, integrating directly with brokers to pull trade data, and
it includes a tax-planning calculator \cite{ref3}. Its breadth of broker
integrations is genuinely useful for the investor segment. The shared limitation
of the category applies here too: it is a closed deterministic product with no
model of notice risk and no natural-language account of its own conclusions. A
user is told what to file, not why, and is not warned about how their return
compares to what the department already holds.

\subsection{Differentiation of Our Approach}
FinITR-AI keeps the one capability this category genuinely does well ---
deterministic, statute-faithful arithmetic --- by embedding its own FY 2025-26
tax engine rather than delegating the numbers to a model. Where it departs is in
two directions that the rule-based tools structurally cannot follow.

\begin{itemize}
    \item It treats AIS-versus-Form-16 reconciliation and quantified notice-risk
          scoring as the primary task, producing an explicit ledger of matched,
          AIS-only, and mismatched items rather than a passive data dump.
    \item It layers genuine reasoning and natural-language explanation over the
          deterministic core, which lets it handle and justify the edge cases a
          fixed rule engine cannot, while still running entirely offline so that
          no document leaves the user's machine.
\end{itemize}

A side-by-side comparison of the category against our approach is given in
Table~\ref{tab:cat-a-comparison}.

\begin{table}[H]
\centering
\caption{Rule-Based / Commercial ITR Tools versus Our Approach}
\label{tab:cat-a-comparison}
\small
\begin{tabularx}{\textwidth}{|l|l|l|X|X|}
\hline
\textbf{System} & \textbf{Type} & \textbf{Year} & \textbf{Methodology} & \textbf{Key Findings / Limitation} \\
\hline
IT e-Filing Portal \cite{ref1} & Govt. utility & 2024 & Schema-validated rule engine with partial AIS pre-fill & Accurate but purely procedural; user must know form and schedules; AIS shown raw \\
\hline
ClearTax \cite{ref2} & Commercial & 2024 & Guided rule-based filing, Form 16 import, regime compare & Shallow AIS reconciliation; closed rules; no notice-risk model; no explanation \\
\hline
Quicko \cite{ref3} & Commercial & 2024 & Guided filing with direct broker capital-gains import & Strong CG import; no risk modelling; closed; no natural-language rationale \\
\hline
\textbf{Proposed (FinITR-AI)} & Agentic + engine & 2026 & Multi-agent reasoning over a deterministic FY25-26 engine, fully offline & AIS reconciliation, quantified notice-risk, citation-verified explanations \\
\hline
\end{tabularx}
\end{table}

\section{Existing Approaches -- General-Purpose Large Language Models}
\label{sec:category-b}

The second category is the direct use of large language models, proprietary or
open, to answer tax and personal-finance questions. The appeal is obvious: no
rigid form, a conversational interface, and an immediate answer. The weakness is
equally structural. These models reason about taxes as text, which means they
inherit two failure modes that matter enormously in this domain --- unreliable
arithmetic, and a tendency to state plausible but incorrect rules with complete
confidence. The distinction worth keeping in mind, and one our evaluation in
Chapter~\ref{chap:results} makes explicit, is that these models often get the
qualitative answer right (which regime, which form) while getting the
quantitative answer wrong (the exact liability), and in a tax filing the
quantitative answer is the one that counts.

\subsection{GPT-4o, Prompted}
GPT-4o is a strong general reasoner, and when prompted with Indian tax scenarios
it frequently identifies the correct direction --- it will usually pick the
right regime and the right form \cite{ref4}. In our own manual evaluation,
however, its failures clustered tightly in the FY 2025-26 specific computations:
it missed the revised marginal-relief threshold, omitted the four percent health
and education cess, and mis-ordered the set-off of short-term capital losses
against gains. These are not careless slips a user can easily audit; they are
confident, well-formatted, wrong final numbers, which is the most hazardous kind
of error in a tool people trust.

\subsection{Gemini 2.0 Flash, Prompted}
Gemini 2.0 Flash exhibits a very similar profile \cite{ref5}: fluent
explanations, correct handling of the well-known rule-recall traps, and the same
class of arithmetic failures on the current year's slabs and surcharge. Its
principal advantage is speed --- its latency is the lowest of the systems we
compared --- which makes it attractive as a conversational assistant. That
advantage does not translate into reliability as a system of record for a tax
computation, which is the role this project cares about.

\subsection{Domain-Specific Financial Language Models}
A separate line of work fine-tunes or instruction-tunes models on financial
corpora. BloombergGPT and the open FinGPT family are the best-known examples
\cite{ref6}. These genuinely improve financial language understanding --- named
entity recognition, sentiment, document classification --- but they are aimed at
markets and reporting rather than at jurisdiction-specific tax computation, and
crucially they still leave the arithmetic to the model rather than to a verified
engine. Specialising the training data does not, on its own, fix the problem
that a neural network is the wrong instrument for applying a tax slab exactly.

\subsection{Differentiation of Our Approach}
FinITR-AI uses a language model, but it is deliberate about confining the model
to what it is good at and removing it from what it is bad at.

\begin{itemize}
    \item \textbf{No model arithmetic.} Every number is computed by a
          deterministic engine; the model never applies a slab or sums a tax.
          The ablation attributes roughly a forty percentage-point gain in tax
          accuracy to this decision alone.
    \item \textbf{Rules pinned to the assessment year.} The FY 2025-26 slabs,
          rebate, surcharge, and marginal relief are hard-coded, so the system
          does not depend on whether a model happens to have the current rules
          embedded in its weights.
    \item \textbf{Active hallucination control.} A Critic agent verifies claims
          against the source documents and a structured tax-law corpus, a step
          that a prompted model, used on its own, does not perform on itself.
\end{itemize}

\begin{table}[H]
\centering
\caption{General-Purpose Large Language Model Approaches}
\label{tab:cat-b-comparison}
\small
\begin{tabularx}{\textwidth}{|l|X|X|}
\hline
\textbf{Model / Method} & \textbf{Key Features} & \textbf{Limitations for Indian ITR} \\
\hline
GPT-4o, prompted \cite{ref4} & Strong general reasoning; good qualitative direction; conversational & FY25-26 slab, cess, and marginal-relief errors; high computational hallucination; cloud-only \\
\hline
Gemini 2.0 Flash, prompted \cite{ref5} & Lowest latency; correct on rule-recall traps & Same arithmetic failures; unreliable as system of record; cloud-only \\
\hline
Domain LLMs (BloombergGPT, FinGPT) \cite{ref6} & Tuned on financial text; stronger finance NLU & Built for markets and reporting, not tax computation; arithmetic still left to the model \\
\hline
\end{tabularx}
\end{table}

\section{Existing Approaches -- Tool-Augmented and Multi-Agent LLM Systems}
\label{sec:category-c}

The third body of work is the research on making language models reliable by
letting them act rather than only speak: calling external tools, retrieving
grounded knowledge, and dividing a task across multiple cooperating agents.
These are the techniques on which FinITR-AI is constructed. The gap they leave
is that they are general scaffolding --- they describe how to reason with tools
and how to ground claims, but they do not supply the tax engine, the
reconciliation logic, or the verification policy that a working filing system
needs.

\subsection{ReAct and Tool-Augmented Reasoning}
The ReAct paradigm interleaves a model's reasoning steps with explicit actions,
so the model can call a tool, read the result, and continue reasoning with that
result in hand \cite{ref7}. Toolformer-style methods extend the idea by teaching
a model when an external call is warranted. This is the direct intellectual
ancestor of our orchestrator, in which the agents reason about what to compute
but always defer the computation itself to a calculator tool. The frameworks
themselves, however, are agnostic to any domain; they say nothing about Indian
tax slabs, about reconciling an AIS against a Form 16, or about when a deduction
must be blocked. Those are the parts this project contributes.

\subsection{Retrieval-Augmented Generation and Faithfulness Verification}
Retrieval-augmented generation grounds a model's output in an external corpus,
retrieving relevant passages and conditioning the generation on them to reduce
hallucination \cite{ref8}. Alongside it, natural-language-inference (NLI) based
faithfulness checking asks whether a generated claim is actually entailed by the
evidence it cites, which provides a mechanical test for grounding \cite{ref9}.
FinITR-AI adopts both ideas concretely: a PageIndex retriever returns the
authoritative text of an Income-Tax-Act section, and an NLI cross-encoder backs
the Critic agent's verdicts. The prior work establishes these as available
components; the step this project takes is to wire them into a closed loop that
can not only flag a weakly-grounded claim but veto it and force the upstream
agent to produce a better one.

\subsection{Differentiation of Our Approach}
FinITR-AI is, in effect, a domain-specific instantiation and tightening of these
general ideas, with two distinguishing features.

\begin{itemize}
    \item \textbf{A specialised division of labour.} Rather than a single
          tool-using agent, the system uses four agents with distinct
          responsibilities --- audit, optimise, comply, and criticise --- each
          grounded in tax-specific tools and a tax-specific corpus, so that the
          reasoning for each sub-problem is handled by a component built for it.
    \item \textbf{Verification with the authority to act.} The Critic does not
          merely annotate low-confidence output. It blocks an ungrounded claim
          and triggers a bounded re-run of the responsible agent, which closes
          the distance between ``retrieval was available'' and ``the
          hallucination was actually prevented.''
\end{itemize}

\begin{table}[H]
\centering
\caption{Tool-Augmented / Multi-Agent LLM Approaches}
\label{tab:cat-c-comparison}
\small
\begin{tabularx}{\textwidth}{|l|X|X|X|}
\hline
\textbf{Method} & \textbf{Key Features} & \textbf{Methodology} & \textbf{Gap We Fill} \\
\hline
ReAct / tool-use \cite{ref7} & Reasoning interleaved with tool calls & Prompted reason--act loops invoking external tools & Domain-agnostic; supplies no tax engine, reconciliation, or verification policy \\
\hline
RAG + NLI faithfulness \cite{ref8}\cite{ref9} & Grounding and entailment-checking of claims & Retrieve evidence; test entailment of generated claims & Components exist; not assembled into a vetoing loop for Indian ITR \\
\hline
\end{tabularx}
\end{table}

To summarise the landscape: the rule-based tools are accurate but blind to the
AIS and unable to explain themselves; the general models are flexible but
arithmetically unreliable and hostile to privacy; and the agentic and retrieval
literature provides the right building blocks but no finished, verified, offline
tax system. FinITR-AI is positioned in the empty middle of this space. It takes
the arithmetic reliability of a rule engine, the reasoning and explanation of a
language model, the AIS reconciliation that neither chatbots nor most tools
attempt, and a verification loop that keeps the whole assembly honest. How that
assembly is engineered is the subject of the next chapter.
